% Institutional background. States no result. Sources: sources/notes.md sections 1.7, 4, 4.1,
% 10 and 11; claims c-ab1228-*, c-council-cap, c-ab257, c-naics, c-coverage-share,
% c-melo-misclass, c-melo-coverage, c-wiltshire, c-d20-federal-floor, c-r50-mw-changes.
\section{Institutional background}
\label{sec:background}

\subsection{AB 1228 and the \$20 fast-food minimum wage}

California's Assembly Bill 1228 was signed on September 28, 2023 and set a minimum wage of \$20
per hour for covered fast-food restaurant employees from April 1, 2024
\citep{dir2026fastfoodcouncil,dir2024mwffsupplement}. It added sections 1474--1476 to the Labor
Code \citep{dir2024fastfoodfaq}. A ``fast food restaurant'' is covered if it meets three
conditions: it is a limited-service restaurant, where customers order and pay before the items
are consumed and there is limited or no table service; it belongs to a chain of at least 60
establishments nationwide; and more than half of its revenue comes from food and beverages for
immediate consumption. Franchisees are covered, and coffee, ice-cream, boba and donut shops can
be. The law exempts restaurants that bake and sell bread on site, restaurants inside grocery
stores of more than 15,000 square feet whose staff are employed by the grocer, and restaurants
in airports, hotels, large event centres, theme parks, museums and gambling establishments, as
well as certain concession locations \citep{dir2024fastfoodfaq}.

The law created a Fast Food Council that may raise the fast-food minimum annually by the lesser
of 3.5 percent and the change in the consumer price index, not before January 1, 2025
\citep{dir2024fastfoodfaq}. No increase took effect through the end of our sample in 2026Q1.
California's general minimum wage orders for 2025 and 2026 state no fast-food dollar rate; they
name fast-food employers under Labor Code section 1474 as subject to a separate minimum
``effective April 1, 2024'' and refer to supplements \citep{dir2026mw}, and the only fast-food
supplement is the 2024 one, at \$20.00 \citep{dir2024mwffsupplement}. The Labor Commissioner's
minimum-wage page, dated December 2025, lists no later fast-food rate
\citep{dir2025minimumwage}, and the Economic Policy Institute reported on March 23, 2026 that
the Council had not yet taken up an adjustment because its chair resigned in June 2025 and no
new chair had been named \citep{bivens2026adjusting}. For the last eight days of 2026Q1 we have
no dated source; a raise would have required a Council vote, and the Council's meeting list
shows none after February 2025 \citep{dir2026fastfoodcouncil}.

Meanwhile California's general minimum wage rose from \$16.00 in 2024 to \$16.50 in 2025 and
\$16.90 in 2026 \citep{dir2026mw}, so the fast-food premium over the general minimum fell from
\$4.00 to \$3.50 and then \$3.10. In our first post year (2024Q2--2025Q1) the premium was \$4.00
in three quarters and \$3.50 in one; in the second (2025Q2--2026Q1) it was \$3.50 in three
quarters and \$3.10 in one. Because the bite of the law fell between the two years, we compare
the year-two employment effect with the year-one effect only next to the yearly wage effects
(hypothesis H3, Section~\ref{sec:decision}).

\subsection{AB 257, the law AB 1228 replaced}
\label{sec:ab257}

AB 1228 replaced AB 257, the FAST Recovery Act, signed on September 5, 2022, which would have let
a Fast Food Council set a minimum wage of up to \$22 per hour in 2023 for fast-food chains with
100 or more locations nationwide, with annual increases of up to the lesser of 3.5 percent and
inflation \citep{melo2026covered,sosinskiy2026effects}. Opponents submitted referendum
signatures in December 2022; a Sacramento court blocked the law in December 2022 and January
2023, and the referendum qualified for the November 2024 ballot on January 24, 2023, which froze
the law before it took effect \citep{ballotpedia2024ab257,melo2026covered,sosinskiy2026effects}.
The compromise announced in September 2023 enacted AB 1228 instead, and the referendum was
withdrawn on December 29, 2023 \citep{sosinskiy2026effects,ballotpedia2024ab257}. No source
reports a wage set under AB 257 \citep{hamdi2025wage,melo2026covered}.

From 2022Q3 a sectoral fast-food wage of up to \$22 was therefore law, and from January 2023 it
was frozen pending a vote. Employers may have responded to it in either direction, and no study
estimates such a response. The quarters 2022Q3--2023Q3 are thus not free of fast-food wage
policy: our base quarter, 2023Q3, is the last quarter before AB 1228 was signed but lies inside
the AB 257 period, as do the September 2023 base of \citet{clemens2026employment} and the
2022Q1--2023Q3 base of \citet[Table~1, Panel~B]{dube2026labor}. Once the COVID quarters
2020Q1--2021Q4 are set aside, only 2022Q1 and 2022Q2 are post-COVID quarters wholly before AB 257
passed. Our design responds in two ways (Section~\ref{sec:strategy}): the selection of the
primary estimator, the level check of the inference and the power calculations use only data
through 2019Q4, and a robustness check moves the base to 2022Q2, the last quarter before AB 257
was enacted.

\subsection{California's general minimum wage before AB 1228}
\label{sec:ramp}

California began raising its general minimum wage in July 2014 and reached \$15 in 2022, with
many local minimums above the state rate \citep{wiltshire2026minimum}. Over the same years the
binding minimum wage in each of the 20 states of the comparison pool we call D20
(Section~\ref{sec:pools}) was the federal \$7.25: on January 1 of every year from 2014 to 2020,
five of them had no state minimum wage law, two had state rates below the federal rate, one had
\$7.25 for larger employers, and twelve had a state rate of exactly \$7.25
\citep{dolwhd2025history,dolwhd2026state}. Of the 30 units outside California and D20, 27 states
and the District of Columbia changed their state minimum wage between January 1, 2014 and
January 1, 2019; Virginia did not, and the Labor Department's January 1 series cannot date
Nevada's change relative to the end of 2019 \citep{dolwhd2025history}. California's pre-AB 1228
path therefore contains the relative effect of its own general minimum-wage ramp, which matters
for how we select the primary estimator (Sections~\ref{sec:select} and~\ref{sec:threats}).

\subsection{What the law covers and what NAICS 722513 measures}
\label{sec:coverage}

The QCEW reports employment by six-digit industry. NAICS 722513, Limited-Service Restaurants,
covers establishments where patrons order or select items and pay before eating, excluding snack
and nonalcoholic beverage bars (722515, such as coffee and donut shops); 722511, Full-Service
Restaurants, covers establishments where patrons are served while seated and pay after eating
\citep{naics2022}. The industry code and the law's coverage differ in three ways.

First, 722513 includes establishments the law does not cover: chains with fewer than 60
locations, independent restaurants and exempt locations. The QCEW does not identify chain size.
No study measures the covered share of 722513 employment directly. \citet{sosinskiy2026effects}
estimate that about 70 percent of the workers in fast-food restaurants and snack and nonalcoholic
beverage bars combined (722513 and 722515) are covered, and give no share for 722513 alone.
\citet{melo2026covered} report that in 2025 employment at California establishments matched by
name to limited-service chains with at least 60 locations, in the modelled establishment data
YTS, is 33 percent smaller than QCEW 722513 employment, which puts covered-chain employment at
roughly two-thirds of the industry code; that comparison comes from a different universe, does
not apply the law's exemptions, and is measured after the law.

Second, covered chains are also coded outside 722513. In Louisiana QCEW microdata matched by
name to fast-food chains, \citet{melo2026covered} find that 77.7 percent of chain restaurant
locations are coded 722513, 9.2 percent 722515, 7.9 percent full-service 722511 and 0.7 percent
722514 or 722320, and that 17 of 121 chains have locations in both 722513 and another restaurant
code. They warn that if classification changes with the composition changes a higher minimum
wage causes, the resulting bias can be more complex than attenuation. We never use 722515, and
722511 is neither a control industry nor a clean placebo in our design
(Section~\ref{sec:robust-plan}).

Third, \citet{melo2026covered} argue that QCEW 722513 may track covered-chain employment poorly
over time and across states: the 2019--2025 percentage change in QCEW 722513 employment minus the
change in their matched-chain YTS employment ranges from $-8$ points (Alabama) to $+17$ points
(District of Columbia), which they note can bias difference-in-differences designs. Because YTS
is itself modelled data, we treat this as a contested objection, not a known bias. It cannot be
removed by our design, and it defines what our estimand is about: an intent-to-treat effect on
QCEW 722513 employment, covered and uncovered establishments together, not the effect on covered
establishments (Section~\ref{sec:estimand}).
